\subsection{A Laplace bound for uniform-matching collisions}
\label{app:lower-proof-matching-laplace}

The next lemma is a finite combinatorial statement independent of the hard
oracle and the adaptive decision rule. It gives a conditional
Laplace bound used below to derive an adaptive lower-tail estimate.

\paragraph{Intuition.}
For a fixed set of $k$ vertices in a uniform matching on $n$ vertices,
each unordered pair is an edge with probability $1/(n-1)$.  Hence the
expected number of internal edges is
\[
 \binom{k}{2}\frac1{n-1}
 =\frac{k(k-1)}{2(n-1)}
 \simeq\frac{k^2}{n}\qquad(k\ge2).
\]
To obtain a lower-tail estimate rather than only a mean, the proof exposes
order $k$ edges, each with conditional success probability at least order
$k/n$.  Their Laplace factors multiply by successive conditioning; the
exposures need not be independent.

\begin{lemma}[Internal-edge Laplace transform]
\label{app:lem:matching-collision-laplace}
Let \(M\) be a uniformly random perfect matching on a labeled set \(U\) of
even cardinality \(n\ge 2\). Fix \(S\subseteq U\), put \(k:=|S|\), and let
\[
 \InternalEdgeCount{S}:=|\{e\in M:e\subseteq S\}|.
\]
Then, for every \(t>0\),
\begin{equation}
 \mathbb E e^{-t\InternalEdgeCount{S}}
 \le
 \exp\left\{
  -\frac{1-e^{-t}}{32}\frac{\PositivePart{k-1}^2}{n}
 \right\}.
 \label{app:eq:matching-collision-laplace}
\end{equation}
\end{lemma}

\begin{proof}

\emph{1. Handle sets of size at most three.}
For \(k\in\{0,1\}\), both sides of
\eqref{app:eq:matching-collision-laplace} equal one. For
\(k\in\{2,3\}\), two internal matching edges would require four distinct
vertices of \(S\). Thus \(\InternalEdgeCount{S}\in\{0,1\}\), and
\[
 \mathbb E\InternalEdgeCount{S}
 =0\cdot\mathbb P(\InternalEdgeCount{S}=0)
  +1\cdot\mathbb P(\InternalEdgeCount{S}=1)
 =\mathbb P(\InternalEdgeCount{S}=1).
\]
Also, writing the edge count as a sum of indicators and taking expectations,
\begin{equation}
 \mathbb E \InternalEdgeCount{S}
 =\sum_{\{u,v\}\subseteq S}\mathbb P(\{u,v\}\in M)
 =\frac{\binom{k}{2}}{n-1}
 =\frac{k(k-1)}{2(n-1)}.
 \label{app:eq:small-set-internal-edge-mean}
\end{equation}
Writing \(a_t:=1-e^{-t}\in(0,1)\), the same two possible values give
\begin{align*}
 \mathbb E e^{-t\InternalEdgeCount{S}}
 &=\mathbb P(\InternalEdgeCount{S}=0)
   +e^{-t}\mathbb P(\InternalEdgeCount{S}=1)\\
 &=1-(1-e^{-t})\mathbb P(\InternalEdgeCount{S}=1)
 =1-a_t\mathbb E\InternalEdgeCount{S}.
\end{align*}
Using \(1-v\le e^{-v}\) for \(v\ge0\), we obtain
\[
 \mathbb E e^{-t\InternalEdgeCount{S}}
 =1-a_t\mathbb E \InternalEdgeCount{S}
 \le e^{-a_t\mathbb E \InternalEdgeCount{S}}
 \le
 \exp\left\{-\frac{a_t}{32}
                   \frac{(k-1)^2}{n}\right\},
\]
where the last inequality follows directly from
\eqref{app:eq:small-set-internal-edge-mean} for \(k=2,3\).


\emph{2. Expose edges and bound each conditional success probability.}
It remains to consider \(k\ge 4\). Partition \(S=A\sqcup B\) with
\(|A|=\lfloor k/2\rfloor\) and \(|B|=\lceil k/2\rceil\), and put
\(m:=\lfloor k/4\rfloor\). We reveal edges of the original matching
\(M\) on the full set \(U\), not of a matching restricted to \(S\).
Only the starting endpoint is required to lie in \(A\); its mate can lie
in \(A\), \(B\), or \(U\setminus S\).

Fix an ordering of \(U\). Let \(U_0:=U\), and for \(j=1,\ldots,m\) define
\begin{align*}
 a_j&:=\min(A\cap U_{j-1}),
 &b_j&:=\operatorname{mate}_{M}(a_j),\\
 e_j&:=\{a_j,b_j\}\in M,
 &U_j&:=U_{j-1}\setminus\{a_j,b_j\}.
\end{align*}
Thus \(U_{j-1}\) consists of the vertices whose matching edges have not yet
been revealed; these vertices are already matched in \(M\).
The minimum uses the fixed ordering, so the choice of \(a_j\) depends only
on the earlier revealed edges. It is well-defined because each previous
edge removes at most two vertices from \(A\), and
\[
 |A\cap U_{j-1}|
 \ge\lfloor k/2\rfloor-2(j-1)
 \ge2m-2(j-1)
 \ge2>0
 \qquad(1\le j\le m).
\]

Define the success indicator by
\[
 X_j:=\mathbf1_{\{b_j\in B\cap U_{j-1}\}},
 \qquad
 \{X_j=1\}=\{\operatorname{mate}_{M}(a_j)\in B\cap U_{j-1}\}.
\]
Since every starting endpoint \(a_\ell\) lies in \(A\), each previous
edge \(e_\ell\) removes at most one vertex from \(B\). Consequently,
\[
 |B\cap U_{j-1}|
 \ge\lceil k/2\rceil-(j-1)
 \ge k/2-(m-1)
 \ge k/4,
 \qquad
 |U_{j-1}|=n-2(j-1).
\]
Conditional on \(e_1,\ldots,e_{j-1}\), the remaining matching is uniform
on \(U_{j-1}\): every perfect matching there has exactly one extension
by the already revealed edges. Each possible mate of \(a_j\) has the same
number of matching completions, so
\[
 \mathbb P(b_j=v\mid e_1,\ldots,e_{j-1})
 =\frac1{|U_{j-1}|-1}
 \quad\text{for }v\in U_{j-1}\setminus\{a_j\}.
\]
Summing over the allowable mates in \(B\), which never include
\(a_j\in A\), gives
\begin{align*}
 \mathbb P(X_j=1\mid e_1,\ldots,e_{j-1})
 &=\sum_{v\in B\cap U_{j-1}}
   \mathbb P(b_j=v\mid e_1,\ldots,e_{j-1})\\
 &=\frac{|B\cap U_{j-1}|}{|U_{j-1}|-1}
 \ge\frac{k/4}{n}=\frac{k}{4n}.
\end{align*}
The earlier indicators \(X_1,\ldots,X_{j-1}\) are determined by these
revealed edges. Adding them to the conditioning therefore gives the same
probability:
\begin{equation}
 \mathbb P(X_j=1\mid X_1,\ldots,X_{j-1},
               \text{all preceding exposed edges})
 \ge u:=\frac{k}{4n}.
 \label{app:eq:matching-cross-success}
\end{equation}
Removing both endpoints after each exposure makes \(e_1,\ldots,e_m\)
distinct. If \(X_j=1\), then \(a_j\in A\) and \(b_j\in B\), so
\(e_j\subseteq A\cup B=S\). Therefore
\[
 \sum_{j=1}^mX_j
 =|\{e_j:1\le j\le m,\ X_j=1\}|
 \le|\{e\in M:e\subseteq S\}|=\InternalEdgeCount{S}.
\]

\emph{3. Iterate conditional Laplace factors without independence.}
For one exposure, the Bernoulli identity and the preceding probability
bound give
\begin{align*}
 &\mathbb E[e^{-tX_j}\mid\text{preceding exposed edges}]\\
 &\qquad=1-(1-e^{-t})\Pr(X_j=1\mid\text{preceding exposed edges})\\
 &\qquad\le1-u(1-e^{-t}).
\end{align*}
For \(1\le r\le m\), condition on the first \(r-1\) exposed edges.
The earlier factors are measurable with respect to this record, so
\begin{align*}
 \mathbb E e^{-t\sum_{j=1}^{r}X_j}
 &=\mathbb E\!\left[
   e^{-t\sum_{j=1}^{r-1}X_j}
   \mathbb E(e^{-tX_r}\mid\text{first \(r-1\) exposed edges})\right]\\
 &\le[1-u(1-e^{-t})]\,
      \mathbb E e^{-t\sum_{j=1}^{r-1}X_j}.
\end{align*}
Starting with the empty sum at \(r=0\), whose exponential is one, and
iterating to \(r=m\) gives the power below without an independence
assumption:
\begin{align}
 \mathbb E e^{-t\InternalEdgeCount{S}}
 &\le \mathbb E\exp\left\{-t\sum_{j=1}^mX_j\right\}\notag\\
 &\le\{1-u(1-e^{-t})\}^{m}
 \le\exp\{-mu(1-e^{-t})\}.
 \label{app:eq:matching-exposure-laplace}
\end{align}
Since \(m\ge k/8\) for \(k\ge4\),
\(mu\ge k^2/(32n)\ge(k-1)^2/(32n)\). Substitution in
\eqref{app:eq:matching-exposure-laplace} proves
\eqref{app:eq:matching-collision-laplace}.
\end{proof}
